\subsection{見積り}

\begin{definitionbox}見積り\par\smallskip
見積りとは、まだ正確には分からない量を、分析的に予測することである。規模、費用、期間、品質、利用量、欠陥数などが対象になる。
\end{definitionbox}

ソフトウェア技術者は、見積りのために見積るのではない。意思決定に必要な量がまだ分からないために見積るのである。見積りは予測であるため、必ず不確実性を持つ。重要なのは、完全に正確な見積りではなく、意思決定を誤らせない程度に十分な見積りである。

\subsubsection{専門家判断}

専門家判断は、専門家の経験と判断に基づく見積りである。最も簡単で、いつでも使える。しかし、単独では精度が低くなりやすい。ワイドバンドデルファイやプランニングポーカーのような集団見積りと組み合わせると精度を高めやすい。

\subsubsection{類推}

類推は、過去に実績が分かっている類似事例を基準にし、差分を調整して見積る方法である。新しい対象を理解し、似た事例を探し、差分を列挙し、差分の大きさを見積り、基準事例の実績へ調整を加える。適切な過去事例がある場合には、専門家判断より精度が高くなりやすい。

\subsubsection{分解}

分解は、対象を小さな部分へ分け、各部分を見積り、合計して全体を求める方法である。ボトムアップ見積りとも呼ばれる。細かく見るため納得感は高いが、作業量が大きい。下位見積りが一方向に偏ると、全体見積りも偏る。

\subsubsection{パラメトリック見積り}

パラメトリック見積りは、観測可能な要因と見積り対象の数学的関係を使う方法である。たとえば、床面積から建設費を求めるように、ソフトウェアでは機能規模、画面数、インターフェース数、複雑度などを使って費用や工数を求める。

この方法は、式が検証されていれば精度が高く、説明しやすい。一方で、式を作るには正確な過去データと統計的な検証が必要である。

\subsubsection{複数見積り}

重要な判断では、ひとつの見積りだけに頼らず、複数の技法や複数の見積り者で見積る。結果が収束していれば信頼しやすい。結果が大きく異なる場合は、重要な要因を見落としている可能性がある。差異の原因を探して再見積りすることで、よりよい意思決定につながる。

\par\smallskip
\begin{center}
\centering
\IfFileExists{assets/generated_figures/ch15/fig-ch15-07.png}{\includegraphics[width=0.96\linewidth]{assets/generated_figures/ch15/fig-ch15-07.png}}{\fbox{\parbox[c][48mm][c]{0.9\linewidth}{\centering 画像生成待ち\\見積り技法の比較表}}}
\par\smallskip
\refstepcounter{figure}\label{fig:ch15-07}図 \thefigure: 見積り技法の比較表
\end{center}
図 \thefigure は、見積り技法の比較表を視覚的に整理したものである。
\par\smallskip
